%% ============================================================
%% Introduction to Cognitive Psychology — Week 5 Study Notes
%% ============================================================
\documentclass[11pt,a4paper]{article}

\usepackage[a4paper, left=1.8cm, right=1.8cm, top=1.3cm, bottom=1.8cm]{geometry}
\usepackage{booktabs}
\usepackage{tabularx}
\usepackage{array}
\usepackage{enumitem}
\usepackage{titlesec}
\usepackage{fancyhdr}
\usepackage{fontenc}
\usepackage{inputenc}
\usepackage{microtype}
\usepackage{parskip}
\usepackage{hyperref}
\usepackage{amsmath}
\usepackage{amssymb}

\hypersetup{hidelinks, colorlinks=false}

%% ── Table helpers ─────────────────────────────────────────
\newcolumntype{L}[1]{>{\raggedright\arraybackslash}p{#1}}

%% ── Section headings ──────────────────────────────────────
\titleformat{\section}{\large\bfseries}{}{0em}{}[\titlerule]
\titleformat{\subsection}{\normalsize\bfseries}{}{0em}{}
\titleformat{\subsubsection}{\normalsize\bfseries}{}{0em}{}

\titlespacing{\section}{0pt}{12pt}{4pt}
\titlespacing{\subsection}{0pt}{8pt}{3pt}
\titlespacing{\subsubsection}{0pt}{6pt}{2pt}

%% ── Page style ────────────────────────────────────────────
\pagestyle{fancy}
\fancyhf{}
\renewcommand{\headrulewidth}{0pt}
\fancyfoot[C]{\small Cognitive Psychology --- Week 5 \textbar\ Page \thepage}

%% ── Misc helpers ──────────────────────────────────────────
\newcommand{\bb}[1]{\textbf{#1}}
\setlength{\parskip}{4pt}

%% ═══════════════════════════════════════════════════════════
\begin{document}
\setlength{\arrayrulewidth}{0.4pt}

\begin{center}
\bfseries\LARGE Introduction to Cognitive Psychology\par
\vspace{0.2cm}
\large Assignment 5 Detailed Solutions
\end{center}
\vspace{0.5cm}

%% ════════════════════════════════════════════════════════════
\section*{ASSIGNMENT 5 --- Full Solutions with Explanations}

\subsubsection*{Question 1}
\textit{As classically conceived, long-term memory is held to have all of the following properties EXCEPT:}
\medskip

\begin{tabular}{L{0.3cm} L{14.9cm}}
& A.\ It comprises a permanent or at least semipermanent store \\
& B.\ It has virtually unlimited capacity \\
& \textbf{C.\ It primarily uses acoustic coding \checkmark} \\
& D.\ Information stored within it is not always easily accessible \\
\end{tabular}

\vspace{0.3cm}
\noindent\textbf{Ans:} C — It primarily uses acoustic coding\\[0.2cm]
\textbf{Why this answer:} \\
Long-term memory does \textbf{not} primarily use acoustic (sound-based) coding. Unlike \textbf{short-term memory}, which relies heavily on \textbf{phonological/acoustic coding} (Conrad, 1964), long-term memory predominantly uses \textbf{semantic (meaning-based) coding}. The other three options are all valid properties of LTM: it serves as a relatively permanent store, it has virtually unlimited capacity, and stored information is not always easily accessible (retrieval failures can occur even when the memory exists).
\vspace{0.4cm}

\subsubsection*{Question 2}
\textit{Psychologists believe that the capacity of long-term memory is:}
\medskip

\begin{tabular}{L{0.3cm} L{14.9cm}}
& \textbf{A.\ Unlimited \checkmark} \\
& B.\ 7 ± 2 items \\
& C.\ 18 items \\
& D.\ 5000 items \\
\end{tabular}

\vspace{0.3cm}
\noindent\textbf{Ans:} A — Unlimited\\[0.2cm]
\textbf{Why this answer:} \\
The capacity of long-term memory is considered \textbf{virtually unlimited} — there is no known upper bound on the amount of information that can be stored. This stands in stark contrast to \textbf{short-term memory}, which has a limited capacity of approximately \textbf{7 ± 2 items} (Miller, 1956).
\vspace{0.4cm}

\subsubsection*{Question 3}
\textit{Which of the following are most likely to be confused in long-term memory?}
\medskip

\begin{tabular}{L{0.3cm} L{14.9cm}}
& A.\ The letters P and R \\
& B.\ The letters C and B \\
& C.\ The words "see" and "bee" \\
& \textbf{D.\ The words "big" and "large" \checkmark} \\
\end{tabular}

\vspace{0.3cm}
\noindent\textbf{Ans:} D — The words "big" and "large"\\[0.2cm]
\textbf{Why this answer:} \\
Because long-term memory is organized primarily by \textbf{semantic (meaning-based) coding}, items that share similar meanings are most likely to be confused. "Big" and "large" are \textbf{synonyms} and are therefore easily confused in LTM. Acoustically similar items like "see" and "bee" cause confusion in \textbf{short-term memory} (where acoustic coding dominates) but not in long-term memory.
\vspace{0.4cm}
\newpage
\subsubsection*{Question 4}
\textit{The code in long-term memory is based on:}
\medskip

\begin{tabular}{L{0.3cm} L{14.9cm}}
& A.\ Sound \\
& B.\ Visual imagery \\
& \textbf{C.\ Meaning \checkmark} \\
& D.\ Both sound and visual imagery \\
\end{tabular}

\vspace{0.3cm}
\noindent\textbf{Ans:} C — Meaning\\[0.2cm]
\textbf{Why this answer:} \\
The dominant coding format in long-term memory is \textbf{semantic coding} — information is stored and organized based on its \textbf{meaning} rather than its sound or visual appearance. LTM errors tend to involve \textbf{semantically related items} (e.g., confusing synonyms) rather than acoustically or visually similar items.
\vspace{0.4cm}

\subsubsection*{Question 5}
\textit{Your memory for how to ride a bicycle is an example of \_\_ memory.}
\medskip

\begin{tabular}{L{0.3cm} L{14.9cm}}
& A.\ Explicit \\
& B.\ Implicit \\
& C.\ Declarative \\
& \textbf{D.\ Procedural \checkmark} \\
\end{tabular}

\vspace{0.3cm}
\noindent\textbf{Ans:} D — Procedural\\[0.2cm]
\textbf{Why this answer:} \\
Memory for \textbf{how to perform a skill} — such as riding a bicycle — is \textbf{procedural memory}. Procedural memory is a subtype of \textbf{implicit (nondeclarative) memory} that stores knowledge about motor skills, habits, and procedures. It is acquired through practice and repetition and is typically difficult to verbalize consciously.
\vspace{0.4cm}

\subsubsection*{Question 6}
\textit{A retrieval cue will be most effective when it is highly distinctive or unusual, according to the principle of:}
\medskip

\begin{tabular}{L{0.3cm} L{14.9cm}}
& \textbf{A.\ Cue overload \checkmark} \\
& B.\ Encoding specificity \\
& C.\ Mood dependence \\
& D.\ State dependence \\
\end{tabular}

\vspace{0.3cm}
\noindent\textbf{Ans:} A — Cue overload\\[0.2cm]
\textbf{Why this answer:} \\
The \textbf{cue overload principle} (Watkins \& Watkins, 1975) states that the effectiveness of a retrieval cue decreases as the number of memories associated with that cue increases. A cue that is \textbf{highly distinctive or unusual} is associated with fewer memories, making it a more effective retrieval aid.
\vspace{0.4cm}
\newpage

\subsubsection*{Question 7}
\textit{Context effects and state-dependent learning effects occur:}
\medskip

\begin{tabular}{L{0.3cm} L{14.9cm}}
& \textbf{A.\ For recall tests only \checkmark} \\
& B.\ For recognition tests only \\
& C.\ For both recall and recognition \\
& D.\ For paired associate tests only \\
\end{tabular}

\vspace{0.3cm}
\noindent\textbf{Ans:} A — For recall tests only\\[0.2cm]
\textbf{Why this answer:} \\
\textbf{Context-dependent} and \textbf{state-dependent} memory effects are primarily observed in \textbf{recall tasks} but not in recognition tasks. \textbf{Recall} relies heavily on retrieval cues — and context or internal state at encoding serves as an important cue. \textbf{Recognition} provides the target item itself as a direct cue, making contextual cues less necessary.
\vspace{0.4cm}

\subsubsection*{Question 8}
\textit{Bahrick's study of retention of Spanish vocabulary words showed that large portions of information remained in long-term memory for:}
\medskip

\begin{tabular}{L{0.3cm} L{14.9cm}}
& A.\ Several months \\
& B.\ 1–2 years \\
& C.\ Up to 5 years \\
& \textbf{D.\ Over 50 years \checkmark} \\
\end{tabular}

\vspace{0.3cm}
\noindent\textbf{Ans:} D — Over 50 years\\[0.2cm]
\textbf{Why this answer:} \\
\textbf{Harry Bahrick's} landmark study (1984) on very long-term memory (termed \textbf{"permastore"}) found that while there was significant forgetting during the first 3–6 years, a substantial portion of Spanish vocabulary remained stable and accessible for \textbf{over 50 years} with virtually no additional decline.
\vspace{0.4cm}

\subsubsection*{Question 9}
\textit{Techniques designed to improve memory, often involving the use of visual imagery, are called:}
\medskip

\begin{tabular}{L{0.3cm} L{14.9cm}}
& \textbf{A.\ Mnemonics \checkmark} \\
& B.\ Eidetic \\
& C.\ Iconic \\
& D.\ IQ enhancers \\
\end{tabular}

\vspace{0.3cm}
\noindent\textbf{Ans:} A — Mnemonics\\[0.2cm]
\textbf{Why this answer:} \\
\textbf{Mnemonics} are structured techniques designed to enhance encoding and retrieval of information. Many strategies rely on \textbf{visual imagery} — for example, the \textbf{method of loci}, the \textbf{peg-word system}, and \textbf{keyword methods}. "Eidetic" refers to photographic-like memory, "iconic" refers to visual sensory memory, and "IQ enhancers" is not a recognized psychological term.
\vspace{0.4cm}

\subsubsection*{Question 10}
\textit{Your memory of your first college lecture would be an example of:}
\medskip

\begin{tabular}{L{0.3cm} L{14.9cm}}
& A.\ Semantic memory \\
& \textbf{B.\ Episodic memory \checkmark} \\
& C.\ Implicit memory \\
& D.\ Working memory \\
\end{tabular}

\vspace{0.3cm}
\noindent\textbf{Ans:} B — Episodic memory\\[0.2cm]
\textbf{Why this answer:} \\
\textbf{Episodic memory} (Tulving, 1972) stores \textbf{personally experienced events} tied to a specific \textbf{time, place, and context.} Your memory of your first college lecture is a personal experience with a particular temporal and spatial context — making it a textbook example of episodic memory, distinct from semantic, implicit, or working memory.
\vspace{0.4cm}

\medskip
\subsection*{Assignment Quick Answer Summary}

\begin{center}
\renewcommand{\arraystretch}{1.4}
\begin{tabular}{|L{0.8cm}|L{6.2cm}|L{8.2cm}|}
\hline
\bb{Q\#} & \bb{Topic} & \bb{Correct Answer} \\
\hline
1 & LTM Properties (NOT acoustic) & C — Acoustic coding \\
2 & LTM Capacity & A — Unlimited \\
3 & Semantic Confusion in LTM & D — "big" and "large" \\
4 & LTM Coding Format & C — Meaning \\
5 & Procedural Memory & D — Procedural \\
6 & Cue Overload Principle & A — Cue overload \\
7 & Context/State-Dependent Effects & A — Recall tests only \\
8 & Bahrick's Permastore Study & D — Over 50 years \\
9 & Mnemonics & A — Mnemonics \\
10 & Episodic Memory & B — Episodic memory \\
\hline
\end{tabular}
\end{center}

\medskip
\end{document}
